\section{Holonomy and odd depth-one values}\label{sec:holonomy}

By Lemma~\ref{lem:rational-descent}, a complex categorical value with
nonzero coefficient of $x^{n-1}y$ gives a rational value with the same
nonvanishing property.  We construct these complex values directly
from a flat connection.
This follows the associator approach to odd-weight elements in
\cite[Proposition~6.3]{Drinfeld}; we include the regularization and
compatibility arguments needed here.  The section first constructs
regularized transport (Proposition~\ref{prop:regularized-holonomy}),
then compares the two transport rules through their operator logarithm,
and finally extracts the nonzero coefficient in
Proposition~\ref{prop:odd-depth-one}.

\subsection{Tree endpoints and regularized transport}
Write $\lambda=(2\pi i)^{-1}$ and
\[
 X_I=\{(z_i)_{i\in I}\in\CC^I:z_i\ne z_j\text{ for }i\ne j\}.
\]
On $X_I$ consider the universal logarithmic braid connection
\cite[Section~1.3]{Kohno}
\begin{equation}\label{eq:holonomy-connection}
 \Omega=\lambda\sum_{i<j}t_{ij}\,d\log(z_i-z_j),
 \qquad dg=\Omega g.
\end{equation}
At every fixed weight this is an ordinary differential equation in a
finite-dimensional nilpotent quotient of the enveloping algebra.
Its solution is the finite sum of iterated integrals in that quotient,
with later differentials multiplying on the left.

The connection is flat.  Its coefficients are closed; terms supported
on four different labels commute; and the remaining terms in
$\Omega\wedge\Omega$ vanish by the three-label braid relations and
\[
 d\log(z_i-z_j)\wedge d\log(z_j-z_k)
 +d\log(z_j-z_k)\wedge d\log(z_k-z_i)
 +d\log(z_k-z_i)\wedge d\log(z_i-z_j)=0.
\]
This is Arnold's logarithmic-form relation \cite{Arnold}.
It follows by putting $u=z_i-z_j$, $v=z_j-z_k$ and using
$z_k-z_i=-u-v$.  The usual homotopy variation formula for an ordinary
differential equation now shows that transport is homotopy invariant:
the variation of transport with fixed endpoints is the integral of
its curvature conjugated by partial transports, and is zero here.
Alternatively this formula follows by differentiating the finite
iterated-integral expression in each weight.
Since $\Omega$ is primitive, transport is group-like.  Indeed its
coproduct and the tensor square of transport solve the same equation
with form $\Omega\otimes1+1\otimes\Omega$ and initial value $1$.

We first define the endpoint data for a fixed choice of scales.
For an internal node $d$ of a tree $p\in\mathcal T(I)$ assign a scale
$s_d=\epsilon^{a_d}$, with $a_{\mathrm{root}}=0$ and exponents strictly
increasing along descent.  Anchor the root at $0$.  The left child of a
node has its anchor at that node's anchor, and the right child has its
anchor displaced by $s_d$; repeat down the tree.  Thus the coordinate
of a leaf is explicitly
\begin{equation}\label{eq:holonomy-tree-coordinate}
 z_i(\epsilon)=
 \sum_{d\text{ above }i}
 \mathbf1_{\{i\text{ lies in the right child of }d\}}\epsilon^{a_d}.
\end{equation}
For sufficiently small positive $\epsilon$ these are distinct real
points in the prescribed leaf order.  The curve
$\epsilon\mapsto(z_i(\epsilon))_{i\in I}$ is the \emph{collar} of $p$.
For example, the three-leaf
objects $a=((12)3)$ and $b=(1(23))$ give respectively
$(0,\epsilon^u,1)$ and $(0,1,1+\epsilon^v)$, with $u,v>0$.
Thus the coordinate $(z_2-z_1)/(z_3-z_1)$ approaches $0$ and $1$
from inside the real interval, as will be used for reassociation.
Let
\[
 K_d=\sum_{\substack{i\text{ in the left child of }d\\
              j\text{ in the right child of }d}}t_{ij},
 \qquad
 E_p(\epsilon)=
 \exp\left(\lambda\sum_dK_d\log s_d\right)
\]
for the object $p$.  All $K_d$ of a fixed tree commute.  Disjoint nodes
are immediate from the disjoint-chord relation.  For nested nodes,
any chord within the smaller cluster commutes with the total of the
chords from that cluster to a fixed outside point, since the two
incident terms give $[t_{ij},t_{ik}+t_{jk}]=0$; summing proves the
claim.

For these chosen collar families, a path between limiting endpoints
means a homotopy class represented by an initial outward collar, a
fixed piecewise smooth path between ordinary configurations, and a
terminal inward collar.  Denote their groupoid by $\mathcal P(I)$.
Truncations at different $\epsilon$ are identified by collar segments.
It is an ordinary fundamental groupoid after identifying each limiting
endpoint with a chosen point on its collar.  Different admissible
exponent systems will be compared below, rather than identified in
this definition.

For $\gamma:p\longrightarrow q$, truncate the collars at $\epsilon$
and write $G(\gamma_\epsilon)$ for ordinary transport.  The proposed
regularized rule is
\begin{equation}\label{eq:holonomy-regularized}
 H_+(\gamma)=\lim_{\epsilon\to0^+}
 E_q(\epsilon)^{-1}G(\gamma_\epsilon)E_p(\epsilon).
\end{equation}
We also consider the rule $H_-$ obtained by conjugating the path and
replacing every chord $t_{ij}$ by $-t_{ij}$; its endpoint frame is
$E_p^{-1}$.  Existence of these limits is part of the next proposition.

The path representatives for the two insertion operations are equally
concrete.  For outer insertion, replace an outer coordinate $z_i(s)$
by the cluster $z_i(s)+\epsilon^M w_{ia}$, where the $w_{ia}$ are fixed
along the outer path at each $\epsilon$ and give the coordinates of
the inserted tree.  For inner insertion, fix the outer
anchors and let one such cluster follow a scaled inner path
$\epsilon^M w_{ia}(s)$.  Use nested smaller scales within each tree,
and choose the inserted scale sufficiently small relative to every
outer separation.  These recipes are used only when the resulting
number of labels is at most four.  The proof below establishes that
they define coherent operations on continued path classes, independent
of the sufficiently separated scale choices.

\begin{proposition}[Regularized transport]\label{prop:regularized-holonomy}
For the tree endpoints and frames just defined,
\eqref{eq:holonomy-regularized} has a limit in every weight.
The limit is independent of the admissible scale exponents under
real interpolation of the endpoint collars.  The insertion recipes
define coherent operations on the continued path classes, and the
limit respects relabeling, outer block insertion, and inner insertion.
The same assertions hold for $H_-$.  The two transport rules agree
on all arity-two paths, including paths reversing the two real points.
\end{proposition}

\begin{proof}
\emph{Convergence at each endpoint.}
If $d$ is the lowest common ancestor of $i,j$, formula
\eqref{eq:holonomy-tree-coordinate} gives
\begin{equation}\label{eq:holonomy-collar-difference}
 z_i-z_j=\sigma_{ij}\epsilon^{a_d}(1+r_{ij}(\epsilon)),
 \qquad \sigma_{ij}\in\{1,-1\},
\end{equation}
where $r_{ij}$ is a finite signed sum of powers
$\epsilon^{a_e-a_d}$ with positive exponents.  Choose a positive lower
bound $\eta$ for all these exponent gaps.  Then
\[
 r_{ij}=O(\epsilon^\eta),\qquad
 r_{ij}'=O(\epsilon^{\eta-1}),\qquad
 \Omega=\lambda\sum_d K_d\,d\log s_d+R(\epsilon)\,d\epsilon,
\]
where every coefficient of $R$ is $O(\epsilon^{\eta-1})$.
The sign $\sigma_{ij}$ contributes no logarithmic differential.

Here and below estimates are at a fixed weight cutoff $N$, using any
norm on the finite-dimensional quotient by weights greater than $N$.
Changing from $g$ to $E_p^{-1}g$ removes the singular term exactly,
because the $K_d$ commute.  The remaining form is
$E_p^{-1}R E_p\,d\epsilon$.  In weight at most $N$, conjugation adds
only polynomial powers of $\log\epsilon$, so its norm is bounded by
\begin{equation}\label{eq:holonomy-collar-bound}
 C_N\epsilon^{\eta-1}(1+|\log\epsilon|)^N\,d\epsilon.
\end{equation}
Its integral between $0$ and $\epsilon$ is
$O_N(\epsilon^\eta(1+|\log\epsilon|)^N)$.
The finite iterated-integral formula, or the differential equation in
the truncated algebra, therefore proves convergence and gives the
same bound for the difference between collar transport in these
frames and the identity.

The collar bound proves the limit in
\eqref{eq:holonomy-regularized}; concatenation, homotopy invariance
and group-likeness pass to the limit.  Different choices of middle
representatives of a continued homotopy class give the same answer by
flatness.

\emph{Independence of endpoint scales.}
To check exponent independence, interpolate two systems of exponents
linearly with a parameter $u\in[0,1]$.  The inequalities between parent
and child exponents persist, and their finitely many gaps have a
uniform positive lower bound $\eta$.  On this interpolation at fixed
$\epsilon$,
\[
 |r_{ij}|\leq C\epsilon^\eta,
 \qquad |\partial_u r_{ij}|
 \leq C\epsilon^\eta|\log\epsilon|.
\]
Again the singular part is precisely $E^{-1}dE$ before changing frame.
The remaining form on the interpolation has integral norm at most
$C_N\epsilon^\eta(1+|\log\epsilon|)^{N+1}$ after changing frame.
Its transport tends to the identity.  Thus identifying endpoints by
these real interpolation paths does not change
\eqref{eq:holonomy-regularized}.

\emph{The two insertion identities.}
We have constructed transport on $\mathcal P(I)$, independently of
the endpoint exponents.  It remains to verify the insertion operations
required of a functor to the chord categories in
Section~\ref{sec:category}.  We first compare the connection forms for
each kind of insertion, then check the operations on path classes.
On a truncated
outer path of the above form, there are constants $c,K>0$ such that
all distinct outer coordinates are separated by at least
$c\epsilon^K$.  Their total variations are bounded independently of
$\epsilon$, and the integral norms of their logarithmic differences
are $O(1+|\log\epsilon|)$.  These statements follow from
\eqref{eq:holonomy-collar-difference} on collars and compactness on
the middle path.

Insert fixed configurations at offsets of size $O(\epsilon^M)$,
with their own smaller nested scales, and choose $M>2K$, also larger
than all outer scale exponents.  Write $D=z_i-z_j$ for an outer
difference and $e$ for the constant difference of the two inserted
offsets.  Then
\begin{equation}\label{eq:holonomy-outer-error}
 d\log(D+e)-d\log D
 =-\frac{e\,dD}{D(D+e)},
 \qquad
 \int\left|d\log(D+e)-d\log D\right|
 \leq C\epsilon^{M-2K}.
\end{equation}
Internal differences are constant along the outer path.  Replacing
all cross-cluster forms by the outer forms gives precisely the
connection obtained from \eqref{eq:holonomy-connection} by
$t_{ij}\mapsto t_{B_iB_j}$.  A difference of iterated integrals of
length at most $N$ is expanded by changing one factor at a time.
The error bound \eqref{eq:holonomy-outer-error} and the logarithmic
bounds for the other factors give an error bounded by a positive
power of $\epsilon$ times a power of $1+|\log\epsilon|$.
Multiplication by the two endpoint frames preserves convergence to
zero: their coefficients of weight at most $N$ are polynomials in
$\log\epsilon$.  A single bound sufficient for all these errors is
\begin{equation}\label{eq:holonomy-insertion-bound}
 C_N\epsilon^\gamma(1+|\log\epsilon|)^{3N},\qquad \gamma>0.
\end{equation}

The full endpoint frame factors exactly as the outer frame under the
block substitution, times the internal frames of the fixed inserted
trees.  The latter frames are the same at both endpoints, and every
inner chord commutes with every outer block sum, as checked in
Section~\ref{sec:category}.  They therefore cancel from normalized
transport.  This proves compatibility with outer insertion.

For inner insertion, keep the outer configuration fixed and insert a
varying inner path $\xi$ scaled uniformly by $\epsilon^M$.  Choose
$M>K$, where $c\epsilon^K$ bounds the distance between its anchor and
outside points from below.  Inner coordinate variations are bounded
as above.  Each cross-cluster differential now has integral norm at
most $C\epsilon^{M-K}$; the inner logarithmic differentials are
unchanged by uniform scaling.  The same iterated-integral argument
gives \eqref{eq:holonomy-insertion-bound}.  There is one additional
frame factor.  If $B$ is the inner label set, it is
\begin{equation}\label{eq:holonomy-central-scale}
 \exp\bigl(\lambda M\log\epsilon\,T_B\bigr),
 \qquad T_B=\sum_{\{i,j\}\subset B}t_{ij}
          =\sum_{d\text{ inner}}K_d.
\end{equation}
Indeed scaling multiplies every inner scale by $\epsilon^M$.
The total chord $T_B$ is central in $\mathfrak t_B$, by summing its
three-label relations.  All outer frame factors commute with inner
chords as well.  These common factors at the two endpoints cancel,
leaving the unscaled inner transport embedded with its original
labels.  The argument includes insertion into a single leaf.
For sufficiently small $\epsilon$ all the cabled paths remain in the
configuration space and vary continuously with $\epsilon$; after
collar identification they consequently define a single continued
homotopy class.  A homotopy of representatives has a compact middle
part, so the inserted scales can be chosen sufficiently small uniformly
on it; applying the same insertion throughout gives a homotopy of the
cabled paths.  Increasing the scale exponent also gives a homotopy
through collision-free paths.  Thus the path operations are well
defined on the continued classes and compatible with concatenation,
and the limiting identities are identities for these operations.

\emph{Coherence of the path operations.}
For truncated outer and inner paths, with parameters $s,t\in[0,1]$
and the uniform bounds above, use the homotopy square
\[
 F(s,t)_{ia}=z_i(s)+\epsilon^M w_{ia}(t).
\]
If $|z_i(s)-z_j(s)|\ge c\epsilon^K$ and
$|w_{ia}(t)|\le C$, then distinct clusters remain separated by
$c\epsilon^K-2C\epsilon^M>c\epsilon^K/2$ for sufficiently small
$\epsilon$ and $M>2K$; within a cluster, differences are
$\epsilon^M(w_{ia}(t)-w_{ib}(t))\ne0$.  Thus the two boundary
routes give the same path class.  Nested insertions agree after
flattening, since
$z+\epsilon^M(w+\epsilon^Nv)=z+\epsilon^Mw+\epsilon^{M+N}v$.
The exponent-interpolation estimate identifies other admissible
scale choices with these, with normalized transport tending to the
identity.  The commuting frame factors above therefore identify
these operations exactly.  For a one-leaf outer tree, restoring
root scale one uses a positive real dilation by $\rho$; its transport
is $\exp(\lambda\log\rho\,T_B)$ and cancels from normalized transport
at the two endpoints by centrality of $T_B$.

\emph{The second transport rule and arity two.}
For $H_-$, apply the same construction to $\bar\gamma$ with
$t_{ij}$ replaced by $-t_{ij}$ and endpoint frame $E_p^{-1}$.
All estimates and cancellation identities just proved remain valid;
relabeling is immediate for both rules.  In arity two the root scale
is $1$, and the difference of the two endpoint coordinates is $1$ or
$-1$.  Every path has logarithmic-difference integral $k\pi i$ for
some integer $k$, of the parity determined by the endpoint orders.
Its first holonomy is $\exp(kt_{12}/2)$.  Conjugating the path negates
$k$, and negating the chord negates it once again, so the second
holonomy is the same.  This checks both full turns and half-turns.
\end{proof}

\subsection{Comparison of the completed path categories}
The two transport rules now respect exactly the operations imposed in
Section~\ref{sec:category}.  To compare them by an automorphism of the
chord category, we first show that each becomes an isomorphism after
completion.
Linearize $\mathcal P(I)$ over $\CC$.
At an object, filter its loop group algebra by powers
of the augmentation ideal; transport this filtration to every arrow
space by any chosen reference path.  Changing the reference path
multiplies by a group element and preserves each filtration step.
Write $\widehat{\CC\mathcal P(I)}$ for the resulting completion, with
coproduct $\Delta\gamma=\gamma\otimes\gamma$.

The following associated-graded comparison is the standard pure-braid
holonomy argument; compare \cite[Proposition~3.6]{BarNatan}.
We give the meridian and completion steps for the path categories just
constructed.

\begin{lemma}\label{lem:holonomy-isomorphism}
Each of $H_+$ and $H_-$ extends to an isomorphism of the completed
linear path categories with $\mathcal C(I;\CC)$.  Both preserve
coproducts and all specified operations, and their associated graded
maps are identical.
\end{lemma}

\begin{proof}
Fix an object, let $P=\pi_1(X_I)$ and let $I_P$ be the augmentation
ideal in $\CC[P]$.  Pair-collision meridians normally generate $P$.
To see this directly, fill a loop by a disk in $\CC^I$ and perturb its
interior to meet the collision hyperplanes transversely in finitely
many smooth points, avoiding their complex-codimension-two
intersections.  Removing small disks about these points expresses the
loop as a product of conjugated positive or negative meridians.
Meridians around the same hyperplane are conjugate, since its smooth
part outside the other hyperplanes is path connected.  It follows
that their classes generate $I_P/I_P^2$.  They are also independent:
the winding numbers of the functions $z_i-z_j$ take value $1$ on the
corresponding positive meridian and $0$ on all the others.

Every augmentation associated graded algebra is generated by its
degree-one part, since $I_P^r/I_P^{r+1}$ is spanned by products of $r$
classes from $I_P/I_P^2$.  These meridian classes satisfy the
infinitesimal braid relations.  Near a disjoint pair of collisions
the two local meridians commute.  Near a triple collision, consider
the orbit loop rotating the small three-point cluster simultaneously
through $2\pi$.  It commutes with every local meridian: a circle
action carries any such loop through a homotopy whose two boundary
composites are the two orders of these loops.  Its degree-one class
is the sum of the three pair-meridian classes, because each of the
three differences winds once and all other differences wind zero
times.  The degree-two commutator therefore says
$[t_{ij},t_{ik}+t_{jk}]=0$ (the commutator with $t_{ij}$ itself is
zero).  Transporting these local loops to the basepoint only changes
them by conjugation, which changes their classes by terms in $I_P^2$
and does not affect their degree-two commutators.  Hence there is a
surjective graded algebra map
\begin{equation}\label{eq:holonomy-meridian-surjection}
 f:U(\mathfrak t_I\otimes\CC)
 \longrightarrow\operatorname{gr}_{I_P}\CC[P].
\end{equation}

Residues show that $H_+$ of a positive pair meridian is
$1+t_{ij}+\text{terms of weight at least two}$: the integral of the
corresponding logarithmic form is $2\pi i$.  For $H_-$ both the turn
and the chord change sign, giving the same result.  Endpoint
normalization conjugates loop transport by a series with constant
term $1$ and hence does not change its degree-one term.  Thus both
holonomies are filtered and each induces a graded map
\[
 g:\operatorname{gr}_{I_P}\CC[P]
 \longrightarrow U(\mathfrak t_I\otimes\CC)
 \quad\text{with}\quad g f(t_{ij})=t_{ij}.
\]
Since the chords generate the enveloping algebra, $gf=\id$.
Surjectivity of $f$ now implies that $f$ and $g$ are inverse
isomorphisms.  In particular both associated graded holonomies are
the same inverse of \eqref{eq:holonomy-meridian-surjection}.

An isomorphism on associated graded spaces lifts inductively to an
isomorphism on every finite filtration quotient: lift the constant
term, then correct the error successively in each weight; injectivity
follows by the first nonzero weight.  Passing to inverse limits proves
the completed assertion on loop algebras.  A reference path identifies
each other arrow space with a loop-algebra torsor, so the same assertion
holds on the whole category.  Coproduct preservation follows from
group-likeness, and preservation of insertions and relabeling follows
from Proposition~\ref{prop:regularized-holonomy}, first on paths and
then on their linear completions by continuity.
\end{proof}

We can now compare the two isomorphisms on the same category:
\[
 S=H_-H_+^{-1}:\mathcal C(I;\CC)\longrightarrow\mathcal C(I;\CC).
\]
It fixes the objects, preserves the coproducts and all operations,
and is the identity on arity two.  Its associated graded is the
identity, so $S-\id$ raises weight on every arrow space.  The operator
\begin{equation}\label{eq:holonomy-operator-log}
 \delta=\log S=
 \sum_{j\geq1}\frac{(-1)^{j+1}}j(S-\id)^j
\end{equation}
is well defined weight by weight and is a compatible derivation and
coderivation.  Here is a direct justification of both rules.  Modulo
weights greater than $N$, the binomial polynomial
\[
 S^s=\sum_{j=0}^{N}\binom{s}{j}(S-\id)^j
\]
agrees with the usual power for every nonnegative integer $s$.
The assertions that it preserve composition, coproduct, and every
insertion are polynomial identities in $s$.  They hold at all such
integers and hence identically.  For the coproduct use the tensor
product truncated by total weight $N$.  Differentiating at $s=0$
gives exactly the derivation and coderivation rules for
\eqref{eq:holonomy-operator-log}, as well as compatibility with the
operations.  Taking homogeneous components gives
$\delta_n\in\mathcal D_n(\CC)$ for all $n>0$.

\subsection{A nonzero depth-one value}
The compatible derivations are now available.  We evaluate their
homogeneous components on the reassociation unit arrow and compare
this operator-logarithm value with the holonomy of a real path.

\begin{proposition}\label{prop:odd-depth-one}
For every odd $n\geq3$ there exists $\psi_n\in V_n$ for which the
coefficient of $x^{n-1}y$ is nonzero.  Equivalently, its ordinary
$y$-depth-one part is a nonzero multiple of $\ad_x^{n-1}y$.
\end{proposition}

\begin{proof}
Use the objects $a,b$ and unit arrow $\alpha$ from
Section~\ref{sec:category}.  Let $\beta$ be the real order-preserving
path of reassociation from $a$ to $b$, and put
\[
 \Phi=H_+(\beta),\qquad U=S(\alpha),\qquad
 \theta=S|_{\operatorname{End}(a)}.
\]
We identify $\Phi,U$ with arrow coefficients.  Since an arrow with
coefficient $h$ is $\alpha h_a$, one has exactly
\begin{equation}\label{eq:holonomy-arrow-action}
 S_{ab}(h)=U\theta(h).
\end{equation}
The path $\beta$ is real, so its conjugate is itself and
$H_-(\beta)=\Phi(-t_{ij})$.  Hence
\begin{equation}\label{eq:holonomy-exact-comparison}
 U\theta(\Phi)=\Phi(-t_{ij}).
\end{equation}
Source $x$ and target $y$ are fixed by the arity-two insertion
operations.  Applying $S$ to $\alpha y_a=y_b\alpha$ therefore yields
\begin{equation}\label{eq:holonomy-source-action}
 \theta(x)=x,\qquad \theta(y)=U^{-1}yU.
\end{equation}
The source total chord $T$ is fixed as well.

The linear term of $\Phi$ is zero.  Along a real order-preserving
path the sign of each difference is constant, so its integral is the
logarithm of the ratio of endpoint absolute differences.  Normalizing
by the frames subtracts the logarithms of their lowest-common-ancestor
scales.  By \eqref{eq:holonomy-collar-difference}, each remaining
ratio tends to $1$.  This proves the claim for every chord separately.
As $\theta$ preserves filtration, equation
\eqref{eq:holonomy-exact-comparison} implies that $U$ too has zero
linear term.  Both series are group-like.  Their logarithms are
primitive and have weight at least two, so
\eqref{eq:category-three-split} places these logarithms in the
completed free Lie algebra on $x,y$.  Every such Lie monomial uses
both letters, since the free Lie algebra on one letter has only
weight one.  Consequently $\Phi,U$ depend only on $x,y$ and
\begin{equation}\label{eq:holonomy-exact-xy}
 U\theta(\Phi)=\Phi(-x,-y),\qquad
 \Phi-1,\ U-1\in J,
\end{equation}
where $J$ is the closed two-sided ideal of series having at least one
$y$ in $\CC\langle\!\langle x,y\rangle\!\rangle$.

Conjugation in \eqref{eq:holonomy-source-action} changes $y$ only by
terms in $J^2$, because $U-1\in J$.  Substitution therefore gives
\[
 (\theta-\id)(J^r)\subseteq J^{r+1}\quad(r\geq1),
 \qquad\theta\equiv\id\pmod{J^2}
\]
on the full series algebra (it fixes the pure-$x$ terms).
Reducing \eqref{eq:holonomy-exact-xy} modulo $J^2$ gives
\begin{equation}\label{eq:holonomy-unit-difference}
 U-1=\Phi(-x,-y)-\Phi(x,y)\pmod{J^2}.
\end{equation}
This is a statement about the image of the unit arrow.  To pass to its
infinitesimal value, put $N=S_{ab}-\id$ on this arrow space.  By
\eqref{eq:holonomy-arrow-action}, for $h\in J$,
\[
 N(h)=(U-1)\theta(h)+(\theta(h)-h)\in J^2.
\]
More generally $N(J^r)\subseteq J^{r+1}$ for $r\geq1$.
Since $N(1)=U-1\in J$, all $N^j(1)$ with $j\geq2$ belong to $J^2$.
The operator logarithm consequently gives
\begin{equation}\label{eq:holonomy-log-depth-one}
 \delta(\alpha)=U-1
 =\Phi(-x,-y)-\Phi(x,y)\pmod{J^2}.
\end{equation}

It remains to compute a single convergent coefficient of $\Phi$.
Set the central $T$ to zero and use
$s=(z_2-z_1)/(z_3-z_1)$.  Writing the three differences as a common
factor times $s,s-1,1$ shows that the common logarithmic differential
has coefficient $T$, while the remaining connection is
\begin{equation}\label{eq:holonomy-interval-connection}
 \lambda\left(x\frac{ds}{s}+y\frac{ds}{s-1}\right),
 \qquad 0<s<1.
\end{equation}
The reassociation path runs from $s=0$ to $s=1$ with positive real
tangents.  With the convention $dg=\Omega g$, the coefficient of
$x^{n-1}y$ has $y$ at the earliest integration time.  Thus for $n\geq2$
it is
\begin{equation}\label{eq:holonomy-zeta-integral}
\begin{split}
 [x^{n-1}y]\Phi
 &=\lambda^n
 \int_{0<s_1<\cdots<s_n<1}
 \frac{ds_1}{s_1-1}\frac{ds_2}{s_2}\cdots\frac{ds_n}{s_n}\\
 &=-\frac{\lambda^n}{(n-1)!}
 \int_0^1\frac{(-\log t)^{n-1}}{1-t}\,dt
 =-\lambda^n\sum_{q\geq1}\frac1{q^n}.
\end{split}
\end{equation}
Here $[w]$ denotes the coefficient of the associative word $w$.
To justify convergence and the last equality, expand
$(1-t)^{-1}=\sum_{q\geq0}t^q$ and apply monotone convergence to the
nonnegative integrand; substitution $t=e^{-u}$ gives the individual
integrals $(n-1)!/(q+1)^n$.  The series converges for $n\geq2$.
Regularization does not change this coefficient.  Its frame factors
are pure-$y$ series on the left and pure-$x$ series on the right; a
nonconstant factor would make a word begin with $y$ or end with $x$,
whereas $x^{n-1}y$ does neither.

Negating both arguments multiplies a weight-$n$ word by $(-1)^n$.
Equations~\eqref{eq:holonomy-log-depth-one} and
\eqref{eq:holonomy-zeta-integral} therefore give, for odd $n\geq3$,
\begin{equation}\label{eq:holonomy-nonzero-odd}
 [x^{n-1}y]\delta_n(\alpha)
 =2\lambda^n\sum_{q\geq1}\frac1{q^n}\ne0.
\end{equation}
This only uses positivity of the real series, with no arithmetic
independence assertion.  The functional $[x^{n-1}y]$ is rational on
the finite data of Lemma~\ref{lem:rational-descent}.  Its nonvanishing
on $\delta_n\in\mathcal D_n(\CC)$ implies nonvanishing on some
$\mathcal D_n(\Q)$, and evaluation gives the required $\psi_n\in V_n$.
Finally, Lie elimination, or direct induction on bracketed words with
one $y$, shows that the weight-$n$, depth-one free Lie space is the
line spanned by $\ad_x^{n-1}y$; its $x^{n-1}y$ coefficient is $1$.
\end{proof}
